% preambulo.tex -- paquetes, macros, notación y entornos de la documentación técnica de QF Desk.
% ÚNICO lugar con \newcommand / \NewDocumentCommand / \DeclareMathOperator. Los capítulos NO definen macros:
% si falta un símbolo se escribe en LaTeX explícito, se declara en la sección «Notación» del capítulo y se
% añade su fila a simbolos/capNN.tex (ver README.md).

% ---------------------------------------------------------------------------------------------------
% Idioma y tipografía. `babel` no tiene instalada la opción `spanish`: se usa `\babelprovide` (títulos en
% español; la partición de palabras usa patrones ingleses, ver README.md).
% ---------------------------------------------------------------------------------------------------
\usepackage[T1]{fontenc}
\usepackage{lmodern}
\usepackage[english]{babel}
\babelprovide[import,main]{spanish}
\usepackage{csquotes}
\usepackage[a4paper,top=2.6cm,bottom=2.6cm,left=2.8cm,right=2.8cm]{geometry}
\usepackage{microtype}
\emergencystretch=2em

% ---------------------------------------------------------------------------------------------------
% Matemáticas, tablas, color y números
% ---------------------------------------------------------------------------------------------------
\usepackage{amsmath,amssymb,amsthm}
\usepackage{booktabs,array,longtable}
\usepackage{xcolor}
\usepackage{xspace}
\usepackage{siunitx}
\sisetup{input-decimal-markers={.}, output-decimal-marker={.}, group-digits=integer}

\definecolor{qfazul}{HTML}{1F4E79}
\definecolor{qfverde}{HTML}{2E6B34}
\definecolor{qfcodigo}{HTML}{4A4A4A}
\definecolor{qfambar}{HTML}{8A5A00}
\definecolor{qfrojo}{HTML}{9C1C1C}
\definecolor{qfgris}{HTML}{3C4A5C}

\numberwithin{equation}{chapter}

% ---------------------------------------------------------------------------------------------------
% Bibliografía (natbib, autor-año; estilo plainnat-es.bst en el directorio del documento)
% ---------------------------------------------------------------------------------------------------
\usepackage[authoryear,round]{natbib}

% ---------------------------------------------------------------------------------------------------
% URL, enlaces y referencias cruzadas (orden obligatorio: xurl, hyperref, doi, cleveref)
% ---------------------------------------------------------------------------------------------------
\usepackage{xurl}
\usepackage[unicode,colorlinks=true,linkcolor=qfazul,citecolor=qfverde,urlcolor=qfazul,
            bookmarksnumbered=true,pdflang={es},pdftitle={QF Desk -- Documentación técnica}]{hyperref}
\usepackage{doi}
\usepackage{cleveref}
\crefname{equation}{ecuación}{ecuaciones}
\Crefname{equation}{Ecuación}{Ecuaciones}
\creflabelformat{equation}{(#2#1#3)}
\crefname{chapter}{capítulo}{capítulos}
\Crefname{chapter}{Capítulo}{Capítulos}
\crefname{section}{sección}{secciones}
\Crefname{section}{Sección}{Secciones}
\crefname{subsection}{subsección}{subsecciones}
\Crefname{subsection}{Subsección}{Subsecciones}
\crefname{table}{tabla}{tablas}
\Crefname{table}{Tabla}{Tablas}
\crefname{figure}{figura}{figuras}
\Crefname{figure}{Figura}{Figuras}
\crefname{appendix}{apéndice}{apéndices}
\Crefname{appendix}{Apéndice}{Apéndices}

% ---------------------------------------------------------------------------------------------------
% Recuadros (tcolorbox) y pendientes (todonotes, solo capítulos 14 y 15)
% ---------------------------------------------------------------------------------------------------
\usepackage{tcolorbox}
\tcbuselibrary{breakable}
\usepackage[colorinlistoftodos,textsize=small]{todonotes}
\setuptodonotes{backgroundcolor=qfambar!12!white,linecolor=qfambar,bordercolor=qfambar}

\newtcolorbox{diferencia}{breakable,colback=qfambar!6!white,colframe=qfambar,
  coltitle=white,fonttitle=\bfseries\small,title={Diferencia con la referencia},
  boxrule=0.6pt,arc=1.5pt,left=6pt,right=6pt,top=4pt,bottom=4pt,before skip=10pt,after skip=10pt}
\newtcolorbox{limites}{breakable,colback=qfgris!5!white,colframe=qfgris,
  coltitle=white,fonttitle=\bfseries\small,title={Límites y supuestos},
  boxrule=0.6pt,arc=1.5pt,left=6pt,right=6pt,top=4pt,bottom=4pt,before skip=10pt,after skip=10pt}
\newtcolorbox{aviso}[1][Aviso]{breakable,colback=qfrojo!5!white,colframe=qfrojo,
  coltitle=white,fonttitle=\bfseries\small,title={#1},
  boxrule=0.8pt,arc=1.5pt,left=6pt,right=6pt,top=4pt,bottom=4pt,before skip=10pt,after skip=10pt}

% ---------------------------------------------------------------------------------------------------
% Fragmentos de código cortos (listings); los acentos se declaran con `literate`
% ---------------------------------------------------------------------------------------------------
\usepackage{listings}
\lstset{basicstyle=\ttfamily\small,breaklines=true,breakatwhitespace=false,columns=fullflexible,
  keepspaces=true,frame=single,rulecolor=\color{qfgris!60!white},xleftmargin=4pt,xrightmargin=4pt,
  literate={á}{{\'a}}1 {é}{{\'e}}1 {í}{{\'i}}1 {ó}{{\'o}}1 {ú}{{\'u}}1 {ñ}{{\~n}}1
           {Á}{{\'A}}1 {É}{{\'E}}1 {Í}{{\'I}}1 {Ó}{{\'O}}1 {Ú}{{\'U}}1 {Ñ}{{\~N}}1}

% ---------------------------------------------------------------------------------------------------
% \codigo{ruta}{símbolo}: referencia al código (R4). Imprime «ruta: símbolo» en monoespaciada, con
% posibles cortes de línea tras `/`, `_`, `.`, `:` y `-`. La ruta es relativa a la raíz del repositorio y
% no lleva números de línea. Se escribe con `_` normal (NO `\_`): scripts de verificación leen el .tex.
% ---------------------------------------------------------------------------------------------------
\ExplSyntaxOn
\cs_new_protected:Npn \qf_codigo_cadena:n #1
  {
    \group_begin:
    \ttfamily
    \str_map_inline:nn {#1}
      {
        \str_case:nnF {##1}
          {
            { / } { /\allowbreak }
            { _ } { \char95 \allowbreak }
            { . } { .\allowbreak }
            { : } { :\allowbreak }
            { - } { -\allowbreak }
          }
          { ##1 }
      }
    \group_end:
  }
\NewDocumentCommand \codigo { m m }
  {
    \textcolor{qfcodigo}{ \qf_codigo_cadena:n {#1} \texttt{:}\space \qf_codigo_cadena:n {#2} }
  }
\ExplSyntaxOff

% \simbolosinput{capNN}: lee simbolos/capNN.tex DENTRO de la tabla de símbolos. El \input de LaTeX deja
% tras de sí marcas no expandibles que abren una celda; el primitivo \@@input es expandible (solo lo usa el
% apéndice A).
\makeatletter
\newcommand{\simbolosinput}[1]{\@@input simbolos/#1.tex }
\makeatother

% \verif{id}{valor}: número de un ejemplo numérico verificable. `id` = capNN-nombre, único en todo el
% documento (no se imprime; lo lee verificacion/comun.py). `valor`: decimal con punto, sin separadores de
% miles, con signo y exponente `e` opcionales. Se imprime tal cual con \num.
\NewDocumentCommand \verif { m m }{\num{#2}}

% ---------------------------------------------------------------------------------------------------
% Notación común (apéndice A). Símbolos sin argumento: se usan en modo matemático ($\Nom$, $\K$);
% fuera de él se envuelven con \ensuremath.
% ---------------------------------------------------------------------------------------------------
\newcommand{\DF}[2]{\ensuremath{P_{#1}(#2)}}            % factor de descuento de la curva c del valor (t_0) a T
\newcommand{\DFt}[3]{\ensuremath{P_{#1}(#2,#3)}}        % factor de descuento de t a T
\newcommand{\zr}[2]{\ensuremath{z_{#1}(#2)}}            % tasa cero
\newcommand{\fwd}[3]{\ensuremath{F_{#1}(#2,#3)}}        % tasa forward simple
\newcommand{\yf}[2]{\ensuremath{\tau(#1,#2)}}           % fracción de año (base explícita en el texto)
\newcommand{\Nom}{\ensuremath{N}}                       % nocional
\newcommand{\K}{\ensuremath{K}}                         % tasa fija / precio de ejercicio
\newcommand{\ON}[1]{\ensuremath{r_{#1}}}                % fijación overnight del día hábil j
\newcommand{\Rcomp}[1]{\ensuremath{R_{#1}}}             % tasa compuesta del periodo i
\newcommand{\Ann}{\ensuremath{A}}                       % anualidad (PV01 sin escalar)
\newcommand{\Spar}{\ensuremath{S}}                      % tasa par
\newcommand{\Spot}{\ensuremath{X_0}}                    % spot USDMXN (MXN por USD)
\newcommand{\FxF}[1]{\ensuremath{X(#1)}}                % forward USDMXN a T
\newcommand{\vol}{\ensuremath{\sigma}}                  % volatilidad
\newcommand{\Ncdf}{\ensuremath{\Phi}}                   % normal estándar acumulada
\newcommand{\PV}{\ensuremath{V}}                        % valor presente
\newcommand{\PnL}{\ensuremath{\mathrm{P\&L}}}           % resultado
\newcommand{\DVO}{\ensuremath{\mathrm{DV01}}}           % sensibilidad a 1 pb
\newcommand{\VaR}[1]{\ensuremath{\mathrm{VaR}_{#1}}}    % valor en riesgo
\newcommand{\ES}[1]{\ensuremath{\mathrm{ES}_{#1}}}      % déficit esperado
\newcommand{\hwa}{\ensuremath{a}}                       % reversión a la media de Hull-White
\newcommand{\hwsig}{\ensuremath{\sigma}}                % volatilidad de Hull-White
\newcommand{\hwtheta}{\ensuremath{\theta(t)}}           % deriva dependiente del tiempo de Hull-White
\newcommand{\basis}{\ensuremath{b}}                     % base de cross currency
\newcommand{\E}{\ensuremath{\mathbb{E}}}                % esperanza
\newcommand{\Q}{\ensuremath{\mathbb{Q}}}                % medida
\newcommand{\cTIIE}{\mathrm{TIIE}}                      % etiquetas de curva (subíndices, modo matemático)
\newcommand{\cSOFR}{\mathrm{SOFR}}
\newcommand{\cMXNcol}{\mathrm{MXN|USD}}
\newcommand{\MXN}{\textup{MXN}\xspace}                  % monedas (texto)
\newcommand{\USD}{\textup{USD}\xspace}
\newcommand{\pb}{\textup{pb}\xspace}                    % punto base; se escribe «1~\pb»
